\section*{अध्याय \ref{ch:g6:shapes} --- त्रिभुज और चतुर्भुज}

\begin{solution}{exo:g6:shapes:1}
तरीक़े का पालन करो: $6$ cm का \omterm{def:g6:lines:objects}{रेखाखंड} $[AB]$, $A$ को केंद्र मानकर
$5$ cm \omterm{def:g3:shapes:shapes}{त्रिज्या} का चाप, $B$ को केंद्र मानकर $4$ cm \omterm{def:g3:shapes:shapes}{त्रिज्या} का
चाप; कटान-बिंदु $C$ है।
\end{solution}

\begin{solution}{exo:g6:shapes:2}
$3$ cm का आधार $[EF]$, फिर $E$ और $F$ को केंद्र मानकर $5$-$5$ cm
\omterm{def:g3:shapes:shapes}{त्रिज्या} के दो चाप; उनका कटान $D$ है। दोनों आधार-कोण
$\widehat{DEF}$ और $\widehat{DFE}$ बराबर हैं (हर एक लगभग
$72.5^\circ$) --- \omterm{def:g6:shapes:triangles}{समद्विबाहु त्रिभुज} की सममिति (\cref{ex:g6:symmetry:proves})।
\end{solution}

\begin{solution}{exo:g6:shapes:3}
$4$ cm का \omterm{def:g6:lines:objects}{रेखाखंड}, फिर उसके सिरों को केंद्र मानकर $4$-$4$ cm
\omterm{def:g3:shapes:shapes}{त्रिज्या} के दो चाप। हर कोण $60^\circ$ का है।
\end{solution}

\begin{solution}{exo:g6:shapes:4}
(a) \omterm{def:g6:shapes:triangles}{समबाहु} (इसलिए \omterm{def:g6:shapes:triangles}{समद्विबाहु} भी)। \quad
(b) \omterm{def:g6:shapes:triangles}{समद्विबाहु} ($5 = 5$)। \quad
(c) \omterm{def:g6:shapes:triangles}{समकोण} और \omterm{def:g6:shapes:triangles}{समद्विबाहु} ($45^\circ$ के दो बराबर कोणों के सामने
दो बराबर भुजाएँ)। \quad
(d) \omterm{def:g6:shapes:triangles}{समबाहु} (तीन बराबर कोण)।
\end{solution}

\begin{solution}{exo:g6:shapes:5}
कर्ण \omterm{def:g6:shapes:triangles}{समकोण} के सामने वाली भुजा है, जो हमेशा सबसे लंबी होती है।
$N$ पर \omterm{def:g6:shapes:triangles}{समकोण} वाले $MNP$ में कर्ण $[MP]$ है: वह \omterm{def:g6:shapes:triangles}{समकोण} वाले \omterm{def:g5:solids:def}{शीर्ष}
$N$ को कभी नहीं छूता।
\end{solution}

\begin{solution}{exo:g6:shapes:6}
\emph{1. सही}: वर्ग के चार \omterm{def:g6:shapes:triangles}{समकोण} होते हैं, और आयत की परिभाषा
यही है।

\emph{2. ग़लत}: $5 \times 3$ के आयत की भुजाएँ बराबर नहीं हैं।

\emph{3. सही}: वर्ग की चार बराबर भुजाएँ होती हैं, जो
\omterm{def:g6:shapes:quadrilaterals}{समचतुर्भुज} की परिभाषा है।

\emph{4. सही}: देखो \cref{ex:g6:shapes:recognize}।
\end{solution}

\begin{solution}{exo:g6:shapes:7}
\cref{ex:g6:shapes:program} की तरह लंबों से रचना।
दोनों विकर्ण लगभग $5.8$ cm के निकलते हैं: आयत के विकर्णों की
लंबाई बराबर होती है (\cref{prop:g6:shapes:diagonals})।
\end{solution}

\begin{solution}{exo:g6:shapes:8}
$6$ cm का \omterm{def:g6:lines:objects}{रेखाखंड} खींचो; उसका मध्यबिंदु $O$; $O$ पर \omterm{def:g6:lines:perp}{लंब};
उस पर $O$ से $2$ cm दूर दोनों ओर एक-एक बिंदु। चारों सिरों को
जोड़ने पर \omterm{def:g6:shapes:quadrilaterals}{समचतुर्भुज} बनता है (चारों \omterm{def:g5:solids:def}{शीर्ष} $O$ से \omterm{def:g6:lines:perp}{लंब} रेखाओं पर
$3$ cm और $2$ cm की दूरियों पर हैं, इसलिए सममिति से चारों भुजाएँ
बराबर हैं)।
\end{solution}

\begin{solution}{exo:g6:shapes:9}
एक सही कार्यक्रम:
1. $BC = 4$ cm वाला \omterm{def:g6:lines:objects}{रेखाखंड} $[BC]$ खींचो;
2. $[BC]$ के ऊपर केंद्र $B$ और \omterm{def:g3:shapes:shapes}{त्रिज्या} $6$ cm का चाप खींचो;
3. $[BC]$ के ऊपर केंद्र $C$ और \omterm{def:g3:shapes:shapes}{त्रिज्या} $6$ cm का चाप खींचो;
4. दोनों चापों के कटान-बिंदु का नाम $A$ रखो;
5. \omterm{def:g6:lines:objects}{रेखाखंड} $[AB]$ और $[AC]$ खींचो।
\end{solution}

\begin{solution}{exo:g6:shapes:10}
वृत्त पर $M$ कहीं भी हो, कोण $\widehat{AMB}$ हमेशा
$90^\circ$ का निकलता है: जब भी $[AB]$ \omterm{def:g4:geometry:circle}{व्यास} हो, त्रिभुज $AMB$
$M$ पर \omterm{def:g6:shapes:triangles}{समकोण} होता है। उपपत्ति \cref{ch:g8:cosine} में मिलेगी।
\end{solution}

\begin{solution}{exo:g6:shapes:11}
ऐसे विकर्ण जो एक-दूसरे को मध्यबिंदु पर काटते हों और जिनकी लंबाई
बराबर हो: यह \emph{आयत} की गुण-सूची है
(\cref{prop:g6:shapes:diagonals})। विकर्णों के \omterm{def:g6:lines:perp}{लंब} होने की कोई
बाध्यता नहीं, इसलिए वह \omterm{def:g6:shapes:quadrilaterals}{समचतुर्भुज} या वर्ग होना ज़रूरी नहीं ---
सही उत्तर आयत है (वर्ग भी शर्त पूरी करेगा, पर उसकी बाध्यता नहीं
है)।
\end{solution}

\begin{solution}{exo:g6:shapes:12}
$4 \times 4$ बिंदुओं के साथ $3 \times 3$ जाल-खाने बनते हैं।
जाल के साथ सीधे वर्ग: भुजाएँ $1$, $2$ या $3$ देने पर
$9 + 4 + 1 = 14$ वर्ग। तिरछे वर्ग भी होते हैं: छोटे वाले, जिनकी
भुजाएँ $1 \times 1$ खानों के विकर्ण हैं (जैसे $(1,0)$, $(2,1)$,
$(1,2)$, $(0,1)$) --- ऐसे चार; और बड़े वाले, जिनकी भुजाएँ
$2 \times 1$ आयतों के विकर्ण हैं, जैसे $(1,0)$, $(3,1)$,
$(2,3)$, $(0,2)$ --- ऐसे दो। ऐसे वर्ग की चारों भुजाएँ बराबर हैं
क्योंकि वे एक जैसे आयतों के विकर्ण हैं। कुल:
$14 + 4 + 2 = 20$ वर्ग। (इस स्तर पर तिरछे वर्ग का एक सही
उदाहरण और $14$ सीधे वर्ग गिन लेना ही सफलता है।)
\end{solution}

\begin{solution}{pb:g6:shapes:1}
\textbf{1.} चारों भुजाएँ एक ही नाप की निकलती हैं (लगभग
$4.2$ cm) और चारों कोण \omterm{def:g6:shapes:triangles}{समकोण} हैं: आकार एक \emph{वर्ग} है।
(बराबर विकर्ण \omterm{def:g6:shapes:triangles}{समकोण} देते हैं, जैसे आयत में; \omterm{def:g6:lines:perp}{लंब} विकर्ण बराबर
भुजाएँ देते हैं, जैसे \omterm{def:g6:shapes:quadrilaterals}{समचतुर्भुज} में; दोनों एक साथ वर्ग देते हैं
--- \cref{prop:g6:shapes:diagonals} उलटी पढ़ी हुई।)

\textbf{2.} एक \emph{\omterm{def:g6:shapes:quadrilaterals}{समचतुर्भुज}}: चारों भुजाएँ बराबर (लगभग
$4.7$ cm), पर कोई \omterm{def:g6:shapes:triangles}{समकोण} नहीं --- ठीक
\cref{exo:g6:shapes:8} वाली रचना।

\textbf{3.} चारों कोण \omterm{def:g6:shapes:triangles}{समकोण} निकलते हैं पर भुजाएँ दो अलग-अलग
लंबाइयों की होती हैं: एक \emph{आयत} --- यही
\cref{exo:g6:shapes:11} वाली स्थिति है।

\textbf{4.} नापने पर बराबर लंबाई का कोई दावा नहीं टिकता और कोई
\omterm{def:g6:shapes:triangles}{समकोण} भी नहीं दिखता; पर आमने-सामने की भुजाएँ दो-दो करके बराबर
हैं और \omterm{def:g6:lines:perp}{समांतर} लगती हैं। यह वही आम ``तिरछा आयत'' है जिसके विकर्ण
बस एक-दूसरे को समद्विभाजित करते हैं --- \emph{\omterm{def:g6:lines:perp}{समांतर} \omterm{def:g4:geometry:polygon}{चतुर्भुज}},
जिसका नाम और अध्ययन \cref{ch:g7:central} में है।

\textbf{5.} विकर्ण अपने साझा मध्यबिंदु पर कटते हुए:
\begin{center}
\renewcommand{\arraystretch}{1.2}
\begin{tabular}{l|l}
विकर्ण\dots & आकार \\
\hline
बराबर और \omterm{def:g6:lines:perp}{लंब} & वर्ग \\
केवल \omterm{def:g6:lines:perp}{लंब} & \omterm{def:g6:shapes:quadrilaterals}{समचतुर्भुज} \\
केवल बराबर & आयत \\
कोई नहीं & बिना नाम वाला (देखो \cref{ch:g7:central}) \\
\end{tabular}
\end{center}

\textbf{6.} कोने के हर त्रिभुज की एक भुजा एक विकर्ण पर है:
$8$ cm का \omterm{def:g3:division:half}{आधा} और $5$ cm का \omterm{def:g3:division:half}{आधा}, यानी $4$ cm और $2.5$ cm की
भुजाएँ, और उनके बीच \omterm{def:g6:shapes:triangles}{समकोण} --- चारों के लिए वही आँकड़े। एक ही दो
भुजाओं और उनके बीच के \omterm{def:g6:shapes:triangles}{समकोण} से बने त्रिभुज एक-दूसरे की हूबहू
नक़ल होते हैं (उन्हें बनाकर देखो: परकार के खुलाव एक जैसे हैं)।
इसलिए उनके चारों कर्ण --- यानी आकार की भुजाएँ --- बराबर हैं:
आकार एक \omterm{def:g6:shapes:quadrilaterals}{समचतुर्भुज} है।

\textbf{7.} परकार को एक भुजा की लंबाई तक खोलो, फिर उसे आकार के
चारों ओर चलाओ: अगर उसकी नोक और पेंसिल हर पड़ोसी कोनों की जोड़ी पर
ठीक बैठें, तो चारों भुजाएँ उसी खुलाव के बराबर हैं। मापपट्टी की
ज़रूरत नहीं।

\textbf{8.} कार्यक्रम:
\begin{enumerate}
  \item $AC = 6$ cm वाला \omterm{def:g6:lines:objects}{रेखाखंड} $[AC]$ खींचो और उसका
        मध्यबिंदु $O$ रखो ($A$ से $3$ cm नापो)।
  \item $O$ पर $(AC)$ पर \omterm{def:g6:lines:perp}{लंब} खींचो।
  \item उस पर $O$ से $3$ cm दूर $B$ और $D$ रखो, हर ओर एक।
  \item $A$, $B$, $C$, $D$ जोड़ दो।
\end{enumerate}
विकर्ण बराबर ($6$ cm), \omterm{def:g6:lines:perp}{लंब}, और एक-दूसरे को समद्विभाजित करते हुए:
कोशिका के अनुसार $ABCD$ एक वर्ग है।

\textbf{9.} नापने पर लगभग $4.2$ cm मिलता है: $3$ cm से
\emph{लंबा} --- आम अंदाज़ा ``विकर्ण का \omterm{def:g3:division:half}{आधा}'' ग़लत है। सिद्ध
किया जा सकने वाला हिस्सा: कोने $A$ से अगले कोने $B$ तक जाने के
लिए या तो भुजा $[AB]$ सीधे चलो, या केंद्र होकर
$A \to O \to B$ का चक्कर लो, जिसमें $3 + 3 = 6$ cm लगते हैं;
सीधी भुजा छोटी है, इसलिए भुजा $6$ cm से कम है।
वह $3$ cm से \emph{ज़्यादा} भी है, यह मापपट्टी दिखा देती है;
\cref{ch:g8:pythagoras} इसे सिद्ध करेगा और ठीक मान भी देगा
($3$ cm गुणा $2$ का वर्गमूल, लगभग $4.24$ cm)।

\textbf{10.} $H$ पर कोनों के चारों त्रिभुज अब \emph{दो} मिलती-जुलती
जोड़ियों में आते हैं: दो में भुजाएँ $3$ cm और $2.5$ cm
(कर्ण $[AB]$ और $[AD]$), दो में भुजाएँ $5$ cm और
$2.5$ cm (कर्ण $[CB]$ और $[CD]$)। तो $AB = AD$ और
$CB = CD$: \omterm{pb:g6:shapes:1}{पतंग} में \emph{पड़ोसी} भुजाओं की दो बराबर जोड़ियाँ
होती हैं (\omterm{def:g6:shapes:quadrilaterals}{समचतुर्भुज} वह ख़ास \omterm{pb:g6:shapes:1}{पतंग} है जिसमें चारों मिल जाती हैं)।

\textbf{11.} \omterm{def:g4:geometry:polygon}{चतुर्भुज}: $2$ विकर्ण। \omterm{def:g4:geometry:polygon}{पंचभुज}: $5$।
\omterm{def:g4:geometry:polygon}{षट्भुज}: $9$।

\textbf{12.} \omterm{def:g4:geometry:polygon}{षट्भुज} के एक \omterm{def:g5:solids:def}{शीर्ष} से विकर्ण ख़ुद उसे और उसके दोनों
पड़ोसियों को छोड़कर हर \omterm{def:g5:solids:def}{शीर्ष} तक जाते हैं: $6 - 3 = 3$
विकर्ण। छह \omterm{def:g5:solids:def}{शीर्ष} मिलकर $6 \times 3 = 18$ विकर्ण-सिरे देते हैं
--- पर हर विकर्ण अपने \emph{दोनों} सिरों पर एक-एक बार गिना गया
है, ठीक वैसे ही जैसे हाथ मिलाना दोनों मिलाने वालों से गिना जाता
है (\cref{pb:g6:wholes:1})। असली गिनती:
$18 \div 2 = 9$। यह चित्र से मेल खाती है।

\textbf{13.} दस \omterm{def:g5:solids:def}{शीर्ष}: हर एक से $10 - 3 = 7$, इसलिए
$10 \times 7 \div 2 = 35$ विकर्ण। बीस \omterm{def:g5:solids:def}{शीर्ष}:
$20 \times 17 \div 2 = 170$ विकर्ण।

\textbf{14.} उम्मीदवार आज़माओ: $10$ \omterm{def:g5:solids:def}{शीर्ष} $35$ देते हैं (प्रश्न
13); $11$ \omterm{def:g5:solids:def}{शीर्ष} $11 \times 8 \div 2 = 44$ देते हैं। मिल गया:
\omterm{def:g4:geometry:polygon}{बहुभुज} के $11$ \omterm{def:g5:solids:def}{शीर्ष} हैं।

\textbf{15.} त्रिभुज का कोई विकर्ण होता ही नहीं: हर \omterm{def:g5:solids:def}{शीर्ष} के
``ग़ैर-पड़ोसी'' हैं\,\dots\ कोई नहीं, क्योंकि बाक़ी दोनों
\omterm{def:g5:solids:def}{शीर्ष} उसके पड़ोसी ही हैं। तरीक़ा भी यही कहता है: हर \omterm{def:g5:solids:def}{शीर्ष} से
$3 - 3 = 0$, कुल $0$। हाथ मिलाना और विकर्ण अलग कपड़ों में एक ही
गिनती हैं: किसी समूह में से चुने गए जोड़े, हर जोड़ा एक बार गिना
हुआ --- गणित अपने अच्छे विचार दोबारा इस्तेमाल करता है।

\end{solution}
